%==========================================================================
% TEMPLATE DE TRABALHO DE CONCLUSÃO DE CURSO — GRADUAÇÃO / FACOM / UFMS
% Baseado na classe abnTeX2
%
% COMO USAR:
%   1. Preencha os campos marcados com <...> na seção "DADOS DO TRABALHO".
%   2. Escreva os capítulos nos arquivos capN/*.tex.
%   3. Adicione suas referências em referencias.bib.
%   4. Compile na ordem: pdflatex -> bibtex -> pdflatex -> pdflatex
%      (ou simplesmente: latexmk -pdf tcc.tex)
%==========================================================================

\documentclass[
    % a opção draft remove as figuras e deixa a compilação mais leve
    %draft,
	% -- opções da classe memoir --
	12pt,				% tamanho da fonte
	openright,			% capítulos começam em pág ímpar (só tem efeito com twoside)
	oneside,			% impressão apenas no anverso. Use twoside para verso e anverso
	a4paper,			% tamanho do papel.
	% -- opções da classe abntex2 --
	%chapter=TITLE,		% títulos de capítulos convertidos em letras maiúsculas
	%section=TITLE,		% títulos de seções convertidos em letras maiúsculas
	%subsection=TITLE,	% títulos de subseções convertidos em letras maiúsculas
	%subsubsection=TITLE,% títulos de subsubseções convertidos em letras maiúsculas
	% -- opções do pacote babel --
	english,			% idioma adicional para hifenização
	brazil,				% o último idioma é o principal do documento
	]{abntex2}

% ---
% PACOTES
% ---

% ---
% Pacotes fundamentais
% ---
\usepackage{cmap}				% Mapear caracteres especiais no PDF
\usepackage{lmodern}			% Usa a fonte Latin Modern
\usepackage[T1]{fontenc}		% Selecao de codigos de fonte.
\usepackage[utf8]{inputenc}		% Codificacao do documento (conversão automática dos acentos)
\usepackage{lastpage}			% Usado pela Ficha catalográfica
\usepackage{indentfirst}		% Indenta o primeiro parágrafo de cada seção.
\usepackage{graphicx}			% Inclusão de gráficos
\usepackage{amssymb}			% Para o comando \checkmark
\usepackage[table,xcdraw]{xcolor}

% Comandos auxiliares para marcação de revisão
\newcommand{\red}[1]{{\color{red} #1}}     % para inserir texto em vermelho
\newcommand{\ma}[1]{{\color{magenta} #1}}  % para inserir texto em magenta
\newcommand{\blue}[1]{{\color{blue} #1}}   % para inserir texto em azul

\usepackage{xurl}				% Quebra de URLs longas em qualquer ponto
\renewcommand{\UrlFont}{\small}
\usepackage{comment}			% incluir trechos de comentários
\usepackage{xspace}
\usepackage{multirow}
\usepackage{hyperref}
\usepackage{pdfpages}			% Inclusão de PDFs (ex.: ficha catalográfica)
\usepackage{pdflscape}
\usepackage{tabularx,ragged2e}
\usepackage{subcaption}
\usepackage{longtable}
\usepackage{rotating}
\usepackage{csquotes}
\usepackage{placeins}
\usepackage{amsmath}
\usepackage{array}
\renewcommand{\arraystretch}{1.2}

% ---
% Listagens de código
% ---
\usepackage{listings}
\lstset{
    basicstyle=\ttfamily\small,      % Define a fonte básica como "small"
    keywordstyle=\color{blue},
    commentstyle=\color{gray},
    stringstyle=\color{green},
    numbers=left,                    % Números nas linhas à esquerda
    numberstyle=\tiny,               % Tamanho da fonte dos números
    stepnumber=1,                    % Intervalo dos números
    breaklines=true,                 % Quebra de linha automática
    frame=single                     % Adiciona uma borda em torno do código
}
\usepackage{caption}
\renewcommand{\lstlistingname}{Código}

% --- Garantir legendas em português ---
\addto\captionsbrazil{%
  \renewcommand{\figurename}{Figura}%
  \renewcommand{\tablename}{Tabela}%
  \renewcommand{\bibname}{Referências}%
}

% --- Pacotes de citação carregados cedo para que abntex2abrev.sty use a
%     definição memoir de \newfloat antes de pacotes como algorithm/float
%     redefinirem o comando. NÃO mova estas duas linhas para o fim do preâmbulo. ---
\usepackage[brazilian,hyperpageref]{backref}
\usepackage[alf]{abntex2cite}		% Citações padrão ABNT (autor-data)
%\usepackage[num]{abntex2cite}		% Alternativa: citações numéricas

% --- Diagramas vetoriais em TikZ ---
\usepackage{tikz}
\usetikzlibrary{arrows.meta, positioning, shapes.geometric, calc, fit}

% --- Algoritmos (pseudocódigo em português) ---
\usepackage{algorithm}
\usepackage{algpseudocode}
\floatname{algorithm}{Algoritmo}
\renewcommand{\algorithmicrequire}{\textbf{Entrada:}}
\renewcommand{\algorithmicensure}{\textbf{Saída:}}
\renewcommand{\algorithmicfor}{\textbf{para}}
\renewcommand{\algorithmicforall}{\textbf{para cada}}
\renewcommand{\algorithmicdo}{\textbf{faça}}
\renewcommand{\algorithmicend}{\textbf{fim}}
\renewcommand{\algorithmicif}{\textbf{se}}
\renewcommand{\algorithmicthen}{\textbf{então}}
\renewcommand{\algorithmicelse}{\textbf{senão}}
\renewcommand{\algorithmicwhile}{\textbf{enquanto}}
\renewcommand{\algorithmicreturn}{\textbf{retorna}}
% Nota: \EndFor, \EndIf e \EndWhile compõem automaticamente "fim para",
% "fim se" e "fim enquanto" a partir das macros acima.

% --- Coluna centralizada de largura fixa: use P{3cm} em tabelas ---
\newcolumntype{P}[1]{>{\centering\arraybackslash}p{#1}}

% --- Geração de texto fictício (remova ao escrever o conteúdo real) ---
\usepackage{lipsum}

% Fallback: define \fonte caso a versão do abnTeX2 não o forneça
\providecommand{\fonte}[1]{\begin{flushleft}\footnotesize{Fonte: #1}\end{flushleft}}

% ---
% CONFIGURAÇÕES DE PACOTES
% ---

% ---
% Configurações do pacote backref
% ---
\renewcommand{\backrefpagesname}{Citado na(s) página(s):~}
\renewcommand{\backref}{}
\renewcommand*{\backrefalt}[4]{
	\ifcase #1 %
		Nenhuma citação no texto.%
	\or
		Citado na página #2.%
	\else
		Citado #1 vezes nas páginas #2.%
	\fi}%
% --- Separadores em PT-BR (evita "and" entre páginas) ---
\renewcommand*{\backreftwosep}{ e~}
\renewcommand*{\backreflastsep}{ e~}
% ---


%==========================================================================
% DADOS DO TRABALHO — PREENCHA OS CAMPOS ABAIXO
%==========================================================================
% Curso e grau. Exemplos:
%   Ciência da Computação         / Bacharel em Ciência da Computação
%   Engenharia de Software        / Bacharel em Engenharia de Software
%   Engenharia de Computação      / Bacharel em Engenharia de Computação
%   Sistemas de Informação        / Bacharel em Sistemas de Informação
%   Análise e Desenvolvimento de Sistemas / Tecnólogo em Análise e Desenvolvimento de Sistemas
\newcommand{\curso}{<Nome do Curso>}
\newcommand{\grau}{Bacharel em <Nome do Curso>}

\titulo{<Título do Trabalho>}
\autor{<Nome Completo do Autor>}
\local{Campo Grande, MS}
\data{<Mês>, <Ano>}
% Para orientadora, use o rótulo opcional:
%   \orientador[Orientadora:]{Profa. Dra. <Nome da Orientadora>}
\orientador{Prof. Dr. <Nome do Orientador>}
% Coorientação (opcional) — descomente a linha correspondente:
%\coorientador{Prof. Dr. <Nome do Coorientador>}
%\coorientadora{Profa. Dra. <Nome da Coorientadora>}
\instituicao{%
  Universidade Federal de Mato Grosso do Sul -- UFMS
  \par
  Faculdade de Computação -- FACOM
  \par
  Curso de \curso}
\tipotrabalho{Trabalho de Conclusão de Curso}
% O preâmbulo deve conter a natureza do trabalho, o objetivo (grau
% pretendido) e o nome da instituição
\preambulo{Trabalho de Conclusão de Curso apresentado ao Curso de
\curso{} da Faculdade de Computação da Universidade Federal de
Mato Grosso do Sul, como requisito parcial para obtenção do grau
de \grau.}
%==========================================================================


% ---
% Configurações de aparência do PDF final
% ---

% alterando o aspecto da cor azul
\definecolor{blue}{RGB}{41,5,195}

% informações do PDF
\makeatletter
\hypersetup{
     	%pagebackref=true,
		pdftitle={\@title},
		pdfauthor={\@author},
    	pdfsubject={\imprimirpreambulo},
	    pdfcreator={LaTeX with abnTeX2},
		pdfkeywords={abnt}{latex}{abntex}{abntex2}{trabalho de conclusão de curso},
		colorlinks=true,       		% false: boxed links; true: colored links
    	linkcolor=blue,          	% color of internal links
    	citecolor=blue,        		% color of links to bibliography
    	filecolor=magenta,      	% color of file links
		urlcolor=blue,
		bookmarksdepth=4
}
\makeatother
% ---

% ---
% Espaçamentos entre linhas e parágrafos
% ---

% O tamanho do parágrafo é dado por:
\setlength{\parindent}{1.3cm}

% Controle do espaçamento entre um parágrafo e outro:
\setlength{\parskip}{0.2cm}  % tente também \onelineskip

% ----
% Início do documento
% ----
\begin{document}

% Retira espaço extra obsoleto entre as frases.
\frenchspacing

% ----------------------------------------------------------
% ELEMENTOS PRÉ-TEXTUAIS
% ----------------------------------------------------------
% \pretextual

% ---
% Capa
% ---
\imprimircapa
% ---

% ---
% Folha de rosto
% (o * indica que haverá a ficha catalográfica no verso)
% ---
\imprimirfolhaderosto*

% ---
% Ficha catalográfica
% Se a biblioteca fornecer a ficha em PDF, comente o ambiente abaixo e use:
%\includepdf{ficha_catalografica.pdf}
% ---
\begin{fichacatalografica}
	\sffamily
	\vspace*{\fill}					% Posição vertical
	\begin{center}					% Minipage Centralizado
	\fbox{\begin{minipage}[c][8cm]{13.5cm}		% Largura
	\small
	\imprimirautor
	%Sobrenome, Nome do autor

	\hspace{0.5cm} \imprimirtitulo  / \imprimirautor. --
	\imprimirlocal, \imprimirdata-

	\hspace{0.5cm} \pageref{LastPage} p. : il. (algumas color.) ; 30 cm.\\

	\hspace{0.5cm} \imprimirorientadorRotulo~\imprimirorientador\\

	\hspace{0.5cm}
	\parbox[t]{12.5cm}{\imprimirtipotrabalho~--~\imprimirinstituicao}\\

	\hspace{0.5cm}
		1. <Palavra-chave 1>.
		2. <Palavra-chave 2>.
		I. \imprimirorientador.
		II. Universidade Federal de Mato Grosso do Sul.
		III. Faculdade de Computação.
		IV. \imprimirtitulo\\

	\hspace{8.75cm} CDU <número>\\
	\end{minipage}}
	\end{center}
\end{fichacatalografica}

% ----------------------------------------------------------
% DEDICATÓRIA (opcional)
% ----------------------------------------------------------
%\begin{dedicatoria}
%   \vspace*{\fill}
%   \centering
%   \noindent
%   \textit{<Dedicatória>} \vspace*{\fill}
%\end{dedicatoria}

% ----------------------------------------------------------
% AGRADECIMENTOS (opcional)
% ----------------------------------------------------------
\begin{agradecimentos}

<Escreva aqui os agradecimentos. É costume agradecer à família, ao(à)
orientador(a), aos professores, aos colegas de curso e à instituição.>

<Se o trabalho recebeu apoio financeiro (por exemplo, bolsa de iniciação
científica, tecnológica ou de extensão), mencione a agência ou o programa
de fomento.>

\end{agradecimentos}

% ----------------------------------------------------------
% EPÍGRAFE (opcional)
% ----------------------------------------------------------
%\begin{epigrafe}
%    \vspace*{\fill}
%    \begin{flushright}
%        \textit{``<Texto da epígrafe>''}\\
%        (<Autor da epígrafe>)
%    \end{flushright}
%\end{epigrafe}

% ----------------------------------------------------------
% RESUMO
% ----------------------------------------------------------
\begin{resumo}

\noindent
<Resumo em português, em parágrafo único, entre 150 e 250 palavras. Deve conter
contextualização, problema, objetivo, método, principais resultados e conclusão.>

\noindent
\textbf{Palavras-chave}: <palavra-chave 1>, <palavra-chave 2>, <palavra-chave 3>.

\end{resumo}

% ----------------------------------------------------------
% ABSTRACT (opcional)
% ----------------------------------------------------------
%\begin{resumo}[Abstract]
%
%\noindent
%<Versão do resumo em inglês, com o mesmo conteúdo da versão em português.>
%
%\noindent
%\textbf{Keywords}: <keyword 1>, <keyword 2>, <keyword 3>.
%
%\end{resumo}

% ----------------------------------------------------------
% LISTA DE FIGURAS
% ----------------------------------------------------------
\pdfbookmark[0]{\listfigurename}{lof}
\listoffigures*
\cleardoublepage

% ----------------------------------------------------------
% LISTA DE TABELAS
% ----------------------------------------------------------
\pdfbookmark[0]{\listtablename}{lot}
\listoftables*
\cleardoublepage

% ----------------------------------------------------------
% LISTA DE ABREVIATURAS E SIGLAS
% ----------------------------------------------------------
% Liste em ordem alfabética. Mantenha apenas as siglas usadas no texto.
\begin{siglas}
  \item[ABNT] Associação Brasileira de Normas Técnicas
  \item[FACOM] Faculdade de Computação
  \item[TCC] Trabalho de Conclusão de Curso
  \item[UFMS] Universidade Federal de Mato Grosso do Sul
  \item[<SIGLA>] <Significado da sigla>
\end{siglas}

% ----------------------------------------------------------
% LISTA DE SÍMBOLOS (opcional)
% ----------------------------------------------------------
%\begin{simbolos}
%  \item[$ \Gamma $] Letra grega Gama
%\end{simbolos}

% ----------------------------------------------------------
% SUMÁRIO
% ----------------------------------------------------------
\pdfbookmark[0]{\contentsname}{toc}
\tableofcontents*
\clearpage
% ---


% ----------------------------------------------------------
% ELEMENTOS TEXTUAIS
% ----------------------------------------------------------
\textual

% ----------------------------------------------------------
% ESTRUTURA DO TCC
% Ajuste a quantidade e a ordem dos capítulos conforme seu trabalho.
% ----------------------------------------------------------

%========================================================
%==  CAPÍTULO 1 - INTRODUÇÃO
%========================================================
\chapter[Introdução]{Introdução}\label{ch:introducao}

<Apresente o contexto geral do trabalho, situando o leitor no tema.>

\section{Problema}\label{sec:problema}

<Descreva o problema que o trabalho pretende resolver ou investigar.>

\section{Justificativa}\label{sec:justificativa}

<Explique por que o problema é relevante e o que motiva o trabalho.>

\section{Objetivos}\label{sec:objetivos}

\subsection{Objetivo Geral}\label{subsec:objetivo_geral}

<Enuncie o objetivo geral em uma frase, iniciada por verbo no infinitivo.>

\subsection{Objetivos Específicos}\label{subsec:objetivos_especificos}

\begin{itemize}
    \item <Objetivo específico 1>;
    \item <Objetivo específico 2>;
    \item <Objetivo específico 3>.
\end{itemize}

% Questões de pesquisa (opcional — útil em trabalhos de natureza investigativa).
%\section{Questões de Pesquisa}\label{sec:questoes_pesquisa}
%
%\begin{itemize}
%    \item[\textbf{QP1}] <Enuncie a primeira questão de pesquisa.>
%    \item[\textbf{QP2}] <Enuncie a segunda questão de pesquisa.>
%\end{itemize}

\section{Organização do Trabalho}\label{sec:organizacao}

<Descreva brevemente o conteúdo de cada capítulo subsequente.>


%========================================================
%==  CAPÍTULO 2 - FUNDAMENTAÇÃO TEÓRICA
%========================================================
\chapter[Fundamentação Teórica]{Fundamentação Teórica}\label{ch:fundamentacao}

<Apresente os conceitos, técnicas e tecnologias necessários para a
compreensão do trabalho.>

\section{<Primeiro Conceito>}\label{sec:conceito1}

<Desenvolva o conceito, apoiando-se em referências \cite{exemplo-artigo}.>

% Exemplo de inclusão de figura. Coloque o arquivo em imagens/.
%\begin{figure}[!ht]
%    \centering
%    \includegraphics[width=0.75\textwidth]{imagens/exemplo.png}
%    \caption{<Legenda da figura>.}
%    \label{fig:exemplo}
%    \fonte{<Autoria própria ou referência>.}
%\end{figure}

\section{<Segundo Conceito>}\label{sec:conceito2}

<Desenvolva o conceito.>

% Exemplo de tabela.
%\begin{table}[!ht]
%\centering
%\caption{<Legenda da tabela>.}
%\label{tab:exemplo}
%\begin{tabular}{l|c|r}
%\hline
%\textbf{Coluna A} & \textbf{Coluna B} & \textbf{Coluna C} \\
%\hline
%<valor> & <valor> & <valor> \\
%\hline
%\end{tabular}
%\fonte{<Autoria própria ou referência>.}
%\end{table}

\section{Trabalhos Relacionados}\label{sec:relacionados}

<Apresente os trabalhos (artigos, outros TCCs, ferramentas ou sistemas
existentes) mais próximos ao seu e explicite o que diferencia a sua proposta.>


%========================================================
%==  CAPÍTULO 3 - DESENVOLVIMENTO
%========================================================
% Este capítulo descreve COMO o trabalho foi realizado. Mantenha apenas
% as seções que se aplicam e renomeie o capítulo se necessário
% (p. ex., "Metodologia", "Proposta", "Solução Proposta").
%
% - Trabalhos de DESENVOLVIMENTO (sistema, aplicativo, ferramenta):
%   use as seções Requisitos, Arquitetura, Implementação e Testes.
% - Trabalhos de PESQUISA (experimento, estudo empírico, avaliação):
%   use as seções Metodologia, Coleta de Dados e Análise.
%========================================================
\chapter[Desenvolvimento]{Desenvolvimento}\label{ch:desenvolvimento}

<Apresente a abordagem adotada e as etapas seguidas para atingir os objetivos.>

%--------------------------------------------------------
% Trabalhos de desenvolvimento
%--------------------------------------------------------
\section{Requisitos}\label{sec:requisitos}

<Descreva os requisitos funcionais e não funcionais da solução.>

\section{Arquitetura e Projeto}\label{sec:arquitetura}

<Apresente a arquitetura da solução, os principais componentes e as decisões
de projeto. Diagramas (UML, TikZ etc.) costumam ser úteis aqui.>

\section{Implementação}\label{sec:implementacao}

<Descreva as tecnologias utilizadas e os aspectos relevantes da implementação.
Indique onde o código-fonte está disponível (por exemplo, um repositório Git).>

% Exemplo de listagem de código.
%\begin{lstlisting}[language=Python, caption={<Legenda do código>}, label={lst:exemplo}]
%def soma(a, b):
%    return a + b
%\end{lstlisting}

% Exemplo de algoritmo em pseudocódigo.
%\begin{algorithm}
%\caption{<Nome do algoritmo>}
%\label{alg:exemplo}
%\begin{algorithmic}[1]
%\Require <entrada>
%\Ensure <saída>
%\For{<condição>}
%    \State <passo>
%\EndFor
%\State \Return <resultado>
%\end{algorithmic}
%\end{algorithm}

\section{Testes e Validação}\label{sec:testes}

<Descreva como a solução foi testada ou validada.>

%--------------------------------------------------------
% Trabalhos de pesquisa (descomente e remova as seções acima, se for o caso)
%--------------------------------------------------------
%\section{Metodologia}\label{sec:metodologia}
%
%<Classifique a pesquisa e descreva o desenho do estudo.>
%
%\section{Coleta de Dados}\label{sec:coleta}
%
%<Descreva o procedimento de coleta de forma reprodutível.>
%
%\section{Análise dos Dados}\label{sec:analise}
%
%<Descreva as técnicas de análise empregadas.>


%========================================================
%==  CAPÍTULO 4 - RESULTADOS E DISCUSSÃO
%========================================================
\chapter[Resultados e Discussão]{Resultados e Discussão}\label{ch:resultados}

<Apresente os resultados obtidos: a solução desenvolvida em funcionamento,
os resultados dos testes ou experimentos, ou os dados analisados.>

\section{<Resultado 1>}\label{sec:resultado1}

<Apresente o resultado com o apoio de figuras e tabelas.>

\section{<Resultado 2>}\label{sec:resultado2}

<Apresente o resultado.>

\section{Discussão}\label{sec:discussao}

<Interprete os resultados, relacionando-os aos objetivos do trabalho e aos
trabalhos relacionados.>


%========================================================
%==  CAPÍTULO 5 - CONCLUSÃO
%========================================================
\chapter[Conclusão]{Conclusão}\label{ch:conclusao}

<Retome o objetivo do trabalho e sintetize o que foi alcançado, indicando
como cada objetivo específico foi atendido.>

\section{Limitações}\label{sec:limitacoes}

<Explicite as limitações do trabalho.>

\section{Trabalhos Futuros}\label{sec:trabalhos_futuros}

<Aponte possíveis continuidades e melhorias.>


% ---
% Finaliza a parte no bookmark do PDF, para que se inicie o bookmark na raiz
% ---
\bookmarksetup{startatroot}%
% ---

% ----------------------------------------------------------
% ELEMENTOS PÓS-TEXTUAIS
% ----------------------------------------------------------
\postextual

% ----------------------------------------------------------
% Referências bibliográficas
% ----------------------------------------------------------
\bibliography{referencias}

% ----------------------------------------------------------
% Apêndices (material elaborado pelo autor) — opcional
% ----------------------------------------------------------
\begin{apendicesenv}

% Imprime uma página indicando o início dos apêndices
\partapendices

%========================================================
%==  APÊNDICE A
%========================================================
\chapter{<Título do Apêndice>}
\label{apx:apendice-a}

<Apêndices contêm material elaborado pelo próprio autor: tabelas extensas,
listagens de dados, instrumentos de coleta, código-fonte etc.>

% Cada \include de um novo apêndice deve ser adicionado ao arquivo principal,
% dentro do ambiente apendicesenv.


\end{apendicesenv}
% ---

% ----------------------------------------------------------
% Anexos (material NÃO elaborado pelo autor) — opcional
% ----------------------------------------------------------
%\begin{anexosenv}
%
%% Imprime uma página indicando o início dos anexos
%\partanexos
%
%\chapter{<Título do Anexo>}
%
%\end{anexosenv}

\end{document}
